\subsection{DOUBLE INVERSE LIMIT}

Let I, L be two preordered sets, and \(I \times L\) their product (§1, no. 4). Consider an inverse system of sets \((\mathbf{E}_{\alpha}^{\lambda}, f_{\alpha\beta}^{\lambda\mu})\) relative to the index set \(I \times L\) .

We have

\[
f _ {\alpha \gamma} ^ {\lambda \nu} = f _ {\alpha \beta} ^ {\lambda \mu} \circ f _ {\beta \gamma} ^ {\mu \nu} \quad \text { whenever } \alpha \leqslant \beta \leqslant \gamma \text { and } \lambda \leqslant \mu \leqslant \nu . \tag {11}
\]

Let \(E\) or \(\lim E_{\alpha}^{\lambda}\) denote the inverse limit of this inverse system.

\[
\overleftarrow {\alpha , \lambda}
\]

For each \(\stackrel{\alpha,}{\lambda} \in \mathbf{L}\) put \(g_{\alpha \beta}^{\lambda} = f_{\alpha \beta}^{\lambda}: \mathbf{E}_{\beta}^{\lambda} \to \mathbf{E}_{\alpha}^{\lambda}\) . It follows from (11) that

\[
g _ {\alpha \gamma} ^ {\lambda} = g _ {\alpha \beta} ^ {\lambda} \circ g _ {\beta \gamma} ^ {\lambda} \quad \text { whenever } \alpha \leqslant \beta \leqslant \gamma ,\tag{12}
\]

so that \((\mathbf{E}_{\alpha}^{\lambda}, g_{\alpha \beta}^{\lambda})\) is an inverse system of sets relative to I. Let \(\mathbf{F}^{\lambda} = \varprojlim_{\alpha} \mathbf{E}_{\alpha}^{\lambda}\) denote its inverse limit. For fixed \(\lambda \leqslant \mu\) in L it follows from (11) that the \(h_{\alpha}^{\lambda \mu} = f_{\alpha \alpha}^{\lambda \mu}: \mathrm{E}_{\alpha}^{\mu} \to \mathrm{E}_{\alpha}^{\lambda}\) form an inverse system of mappings, whose inverse limit we denote by \(h^{\lambda \mu} = \varprojlim_{\alpha} h_{\alpha}^{\lambda \mu}: \mathrm{F}^{\mu} \to \mathrm{F}^{\lambda}\) . If \(\lambda \leqslant \mu \leqslant \nu\) in L, we have

\[
h ^ {\lambda \nu} = h ^ {\lambda \mu} \circ h ^ {\mu \nu}\tag{13}
\]

(no. 2, Proposition 2, Corollary 2); hence \((\mathbf{F}^{\lambda}, h^{\lambda \mu})\) is an inverse system of sets relative to L. Let \(\mathbf{F} = \lim_{\lambda \in \mathbf{L}} \mathbf{F}^{\lambda}\) be its inverse limit. We shall define a canonical bijection \(\mathbf{F} \to \mathbf{E}\) . To do this we note that \(\mathbf{F}\) is by definition a subset of \(\prod_{\lambda \in \mathbf{L}} \mathbf{F}^{\lambda}\) , and \(\mathbf{F}^{\lambda}\) is a subset of \(\prod_{\alpha \in \mathbf{I}} \mathbf{E}_{\alpha}^{\lambda}\) ; hence \(\mathbf{F}\) may be canonically identified with a subset of \(\prod_{(\alpha, \lambda) \in \mathbf{I} \times \mathbf{L}} \mathbf{E}_{\alpha}^{\lambda} = \mathbf{G}\) (Chapter II, §5, no. 5, Proposition 7). For each \(z \in \mathbf{G}\) , let \(\operatorname{pr}^{\lambda}(z)\) denote the element \((\operatorname{pr}_{\alpha}^{\lambda}(z))_{\alpha \in \mathbf{I}}\) of \(\prod_{z \in \mathbf{I}} \mathbf{E}_{\alpha}^{\lambda}\) . Then \(z \in \mathbf{F}\) if and only if

\[
\operatorname{pr} ^ {\lambda} (z) = h ^ {\lambda \mu} (\operatorname{pr} ^ {\mu} (z)) \quad \text { whenever } \lambda \leqslant \mu \text { in } L\tag{14}
\]

and \(\mathbf{pr}^{\lambda}(z) \in \mathbf{F}^{\lambda}\) for all \(\lambda \in \mathbf{L}\) ; that is to say, whenever \(\alpha \leqslant \beta\) in I we have

\[
\operatorname{pr} _ {\alpha} ^ {\lambda} (z) = f _ {\alpha \beta} ^ {\lambda \lambda} (\operatorname{pr} _ {\beta} ^ {\lambda} (z)).\tag{15}
\]

But \(h^{\lambda \mu}(\mathrm{pr}^{\mu}(z)) = (f_{\alpha \alpha}^{\mu \lambda}(\mathrm{pr}_{\alpha}^{\mu}(z))_{\alpha \in \mathbf{I}}\) ; it therefore follows from (14) and (15) that if \(\alpha \leqslant \beta\) and \(\lambda \leqslant \mu\) , we have

\[
\operatorname{pr} _ {\alpha} ^ {\lambda} (z) = f _ {\alpha \alpha} ^ {\lambda \mu} (f _ {\alpha \beta} ^ {\mu \mu} (\operatorname{pr} _ {\beta} ^ {\mu} (z))) = f _ {\alpha \beta} ^ {\lambda \mu} (\operatorname{pr} _ {\beta} ^ {\mu} (z)),
\]

which implies that \(z \in \mathbf{E}\) . The converse is obvious, and we have therefore proved

PROPOSITION 4. If \((\mathbf{E}_{\alpha}^{\lambda}, f_{\alpha \beta}^{\lambda \mu})\) is an inverse system of sets relative to a product \(\mathbf{I} \times \mathbf{L}\) of preordered sets, then (up to a canonical bijection) we have

\[
\varprojlim_ {\alpha , \lambda} \mathrm{E} _ {\alpha} ^ {\lambda} = \varprojlim_ {\lambda} (\varprojlim_ {\alpha} \mathrm{E} _ {\alpha} ^ {\lambda}).\tag{16}
\]

COROLLARY 1. Let \((\mathbf{E}_{\alpha}^{\prime \lambda}, f_{\alpha \beta}^{\prime \mu \lambda})\) be another inverse system of sets relative to \(\mathbf{I} \times \mathbf{L}\) , and for each \((\alpha, \lambda) \in \mathbf{I} \times \mathbf{L}\) let \(u_{\alpha}^{\lambda}\) be a mapping of \(\mathbf{E}_{\alpha}^{\lambda}\) into \(\mathbf{E}_{\alpha}^{\prime \lambda}\) such that the \(u_{\alpha}^{\lambda}\) form an inverse system of mappings. Then

\[
\varprojlim_ {\alpha , \lambda} u _ {\alpha} ^ {\lambda} = \varprojlim_ {\lambda} (\varprojlim_ {\alpha} u _ {\alpha} ^ {\lambda}).\tag{17}
\]

The verification is similar to that of Proposition 4.

COROLLARY 2. Let \((\mathrm{E}_{\alpha}^{\lambda}, f_{\alpha \beta}^{\lambda})_{\lambda \in \mathbf{L}}\) be a family of inverse systems of sets relative to I. If \(\prod_{\lambda \in \mathbf{L}} f_{\alpha \beta}^{\lambda}\) denotes the extension to products (Chapter II, §5, no. 7, Definition 2) of the family of mappings \((f_{\alpha \beta}^{\lambda})_{\lambda \in \mathbf{L}}\) , then \(\left(\prod_{\lambda \in \mathbf{L}} \mathrm{E}_{\alpha}^{\lambda}, \prod_{\lambda \in \mathbf{L}} f_{\alpha \beta}^{\lambda}\right)\) is an inverse system of sets relative to I, and (up to a canonical bijection) we have

\[
\varprojlim_ {\alpha} \prod_ {\lambda \in \mathbf {L}} \mathrm{E} _ {\alpha} ^ {\lambda} = \prod_ {\lambda \in \mathbf {L}} (\varprojlim_ {\alpha} \mathrm{E} _ {\alpha} ^ {\lambda}).\tag{18}
\]

Consider the double inverse system \((\mathbf{E}_{\alpha}^{\lambda}, g_{\alpha\beta}^{\lambda\mu})\) relative to \(I \times L\) , where the order relation on L is equality (no. 1, Example 1), and apply Proposition 4.

